% Contenido del póster: Beaumont, Tang y Vansteenberghe (2026), The Review of Financial Studies. DOI 10.1093/rfs/hhaf080
% Elaborado con apoyo de IA a partir del resumen oficial del artículo (OpenAlex), 2026-11-12.
\def\DOI{10.1093/rfs/hhaf080}
\def\TituloEs{Efectos de garantía: el papel de las fintech en el crédito a pequeñas empresas}
\def\TituloOriginal{Collateral effects: The role of FinTech in small business lending}
\def\Autores{Beaumont, Tang y Vansteenberghe}
\def\Anio{2026}
\def\Revista{The Review of Financial Studies}

\def\Problema{Las fintech otorgan préstamos subordinados y sin garantía a pequeñas empresas. No se sabe si esos préstamos desplazan al crédito bancario o si, por el contrario, lo facilitan.}
\def\Objetivo{Investigar el impacto de introducir préstamos fintech subordinados y sin garantía en el mercado de crédito a pequeñas empresas.}
\def\Metodo{Datos administrativos de Francia. Para controlar la demanda de crédito, los autores usan emparejamiento por puntaje de propensión y un instrumento de tipo \textit{shift-share}.}
\def\Resultados{\begin{itemize}
  \item Tras recibir un préstamo fintech, el crédito bancario de la empresa aumenta 13~\%.
  \item El aumento solo ocurre cuando la empresa invierte en activos nuevos.
  \item Después, los prestatarios fintech tienen más probabilidad de entregar garantías a los bancos.
\end{itemize}}
\def\Conclusion{Las empresas usan el préstamo fintech para comprar activos que luego dan en garantía a los bancos, con lo que aumenta su capacidad total de endeudamiento: el crédito fintech complementa al bancario.}
\def\Relevancia{Los proyectos empresariales del financiamiento participativo peruano recibieron en promedio S/~\cifra{montoMedioEmp} mil. Para saber si, como en Francia, ese financiamiento facilita luego el crédito bancario, haría falta seguir a cada empresa en el sistema financiero, información que no es pública. Es una línea de investigación posible con datos administrativos.}
\def\PalabrasClave{fintech, crédito a pequeñas empresas, garantías, complementariedad, Francia}
% --- Ficha del estudio (herramientas/posteres/ficha_articulos.py) ---
\def\Ambito{Francia (datos administrativos)}
\def\Enfoque{Empírico: emparejamiento por puntaje de propensión e instrumento \textit{shift-share}}
\def\FiguraPeru{2}
% --- fin de la ficha ---
